\question[4]
Gegeben ist das folgende Programm.
\begin{lstlisting}
import tkinter as tk

def klick():
    zahl = int(anzeige['text'])
    anzeige['text'] = zahl * 2

fenster = tk.Tk()
anzeige = tk.Label(fenster, text='3')
anzeige.pack()
knopf = tk.Button(fenster, text='Mal 2', command=klick)
knopf.pack()
fenster.mainloop()
\end{lstlisting}
\begin{parts}
\part Was steht im Fenster, bevor geklickt wird?
\part Was zeigt das Label nach drei Klicks auf den Knopf?
\part Warum steht in der Funktion \texttt{int(anzeige['text'])}?
\end{parts}
\begin{solutionbox}{3cm}
a) das Label mit 3 und der Knopf \texttt{Mal 2} \\
b) 24 \\
c) Das Label enthält Text. Zum Rechnen muss er in eine ganze Zahl umgewandelt werden.
\end{solutionbox}
